% GENERATED FILE. Edit canonical lessons, then run npm run generate:books.
% canonical-source-hash: fnv1a64:c6f8e47bc84bf9d9
% canonical-lessons: MR-C441-ksitij, MR-C441-candne, MR-C441-ugavne, MR-C441-mavalne, MR-C441-camakne

\chapter{Sunrise and Moonlight}
\label{ch:mr-sunrise-and-moonlight}
\hlchaptermodality{441}

\begin{chapteropening}
\textbf{By the end of this chapter:} \emph{I can talk about the horizon and moonlight, and say the sun rose, set and shone.}

Complete the last lesson of chapter 441: I can talk about the horizon and moonlight, and say the sun rose, set and shone.
\end{chapteropening}

\section[\texorpdfstring{\mr{क्षितिज} (kṣitij)}{क्षितिज (kṣitij)}]{\mr{क्षितिज} (kṣitij) — the horizon}

\label{lesson:MR-C441-ksitij}



\begin{quote}
Before the new one: say the Marathi for to lash, then the Marathi for to recede.
\end{quote}

\subsection*{You'll want to know: \mr{क्षितिज}}
\textbf{\mr{क्षितिज}} — \emph{kṣitij} — \textquotedblleft{}the horizon\textquotedblright{}.

\begin{quote}
\mr{समुद्रावर} \mr{क्षितिज} \mr{लांबवर} \mr{दिसतं}.
\end{quote}

\begin{quote}
\emph{samudrāvar kṣitij lāmbvar distaṁ.}
\end{quote}

\begin{quote}
At sea the horizon looks far away.
\end{quote}

Neuter.

A neuter noun: \mr{क्षितिजावर}, on the horizon.

\begin{cousinweb}
\textbf{\mr{क्षितिज}} (\emph{kṣitij}) is a tatsama, taken from Sanskrit \emph{kṣitija}, \textquotedblleft{}born of the earth\textquotedblright{} (Monier-Williams), used in Marathi for the line where earth meets sky.
\end{cousinweb}

\begin{grammarlens}[title={What you've built}]
One more word for getting around.
\end{grammarlens}

\subsection*{Your turn}
\begin{itemize}
\raggedright
  \item \emph{Say it:} \emph{kṣitij}
  \item \emph{Say it:} \emph{kṣitij}, once more
  \item \emph{Say it:} \emph{samudrāvar kṣitij lāmbvar distaṁ.} — At sea the horizon looks far away.
  \item \emph{Recall:} say the Marathi for to lash, then the Marathi for to recede, then say \emph{kṣitij} again
\end{itemize}

\begin{tcolorbox}[breakable,colback=teal!4,colframe=teal!35!black,title={Before you move on}]
What does \mr{क्षितिज} mean? (\textquotedblleft{}The horizon\textquotedblright{}.) Say it once more.
\end{tcolorbox}
\section[\texorpdfstring{\mr{चांदणे} (cāndṇe)}{चांदणे (cāndṇe)}]{\mr{चांदणे} (cāndṇe) — moonlight}

\label{lesson:MR-C441-candne}



\begin{quote}
Before the new one: say the Marathi for to recede, then the Marathi for the horizon.
\end{quote}

\subsection*{You'll want to know: \mr{चांदणे}}
\textbf{\mr{चांदणे}} — \emph{cāndṇe} — \textquotedblleft{}moonlight\textquotedblright{}.

\begin{quote}
\mr{अंगणात} \mr{चांदणं} \mr{पडलं} \mr{होतं}.
\end{quote}

\begin{quote}
\emph{aṅgaṇāt cāndṇaṁ paḍlaṁ hotaṁ.}
\end{quote}

\begin{quote}
Moonlight lay in the courtyard.
\end{quote}

Neuter.

Written \mr{चांदणे}, said and often written \mr{चांदणं}, like other neuters in -\mr{णे}.

\begin{cousinweb}
\textbf{\mr{चांदणे}} (\emph{cāndṇe}) is inherited: CDIAL lists Marathi \emph{ċā̃dṇẽ}, \textquotedblleft{}moonlight\textquotedblright{}, under \emph{*cāndraṇa}, \textquotedblleft{}moonlight\textquotedblright{}, a derivative of Sanskrit \emph{candra}, \textquotedblleft{}the moon\textquotedblright{} (CDIAL 4745).
\end{cousinweb}

\begin{grammarlens}[title={What you've built}]
One more word for getting around.
\end{grammarlens}

\subsection*{Your turn}
\begin{itemize}
\raggedright
  \item \emph{Say it:} \emph{cāndṇe}
  \item \emph{Say it:} \emph{cāndṇe}, once more
  \item \emph{Say it:} \emph{aṅgaṇāt cāndṇaṁ paḍlaṁ hotaṁ.} — Moonlight lay in the courtyard.
  \item \emph{Recall:} say the Marathi for to recede, then the Marathi for the horizon, then say \emph{cāndṇe} again
\end{itemize}

\begin{tcolorbox}[breakable,colback=teal!4,colframe=teal!35!black,title={Before you move on}]
What does \mr{चांदणे} mean? (\textquotedblleft{}Moonlight\textquotedblright{}.) Say it once more.
\end{tcolorbox}
\section[\texorpdfstring{\mr{उगवणे} (ugavṇe)}{उगवणे (ugavṇe)}]{\mr{उगवणे} (ugavṇe) — to rise; to sprout}

\label{lesson:MR-C441-ugavne}



\begin{quote}
Before the new one: say the Marathi for the horizon, then the Marathi for moonlight.
\end{quote}

\subsection*{You'll want to know: \mr{उगवणे}}
\textbf{\mr{उगवणे}} — \emph{ugavṇe} — \textquotedblleft{}to rise; to sprout\textquotedblright{}.

\begin{quote}
\mr{आज} \mr{सूर्य} \mr{सहा} \mr{वाजता} \mr{उगवला}.
\end{quote}

\begin{quote}
\emph{āj sūrya sahā vājtā ugavlā.}
\end{quote}

\begin{quote}
Today the sun rose at six.
\end{quote}

Past (he …): \mr{तो} \mr{उगवला}, \emph{to ugavlā}.

For the sun and moon, and for seeds coming up.

\begin{cousinweb}
\textbf{\mr{उगवणे}} (\emph{ugavṇe}) is inherited: CDIAL lists Marathi \emph{ugavṇẽ}, \textquotedblleft{}to rise (of sun), shoot up\textquotedblright{}, under \emph{*udgamati}, \textquotedblleft{}comes forth, rises\textquotedblright{} (CDIAL 1947).
\end{cousinweb}

\begin{grammarlens}[title={What you've built}]
One more word for getting around.
\end{grammarlens}

\subsection*{Your turn}
\begin{itemize}
\raggedright
  \item \emph{Say it:} \emph{ugavṇe}
  \item \emph{Say it:} \emph{ugavṇe}, once more
  \item \emph{Say it:} \emph{āj sūrya sahā vājtā ugavlā.} — Today the sun rose at six.
  \item \emph{Recall:} say the Marathi for the horizon, then the Marathi for moonlight, then say \emph{ugavṇe} again
\end{itemize}

\begin{tcolorbox}[breakable,colback=teal!4,colframe=teal!35!black,title={Before you move on}]
What does \mr{उगवणे} mean? (\textquotedblleft{}To rise; to sprout\textquotedblright{}.) Say it once more.
\end{tcolorbox}
\section[\texorpdfstring{\mr{मावळणे} (māvaḷṇe)}{मावळणे (māvaḷṇe)}]{\mr{मावळणे} (māvaḷṇe) — to set}

\label{lesson:MR-C441-mavalne}



\begin{quote}
Before the new one: say the Marathi for moonlight, then the Marathi for to rise.
\end{quote}

\subsection*{You'll want to know: \mr{मावळणे}}
\textbf{\mr{मावळणे}} — \emph{māvaḷṇe} — \textquotedblleft{}to set\textquotedblright{}.

\begin{quote}
\mr{सूर्य} \mr{मावळला} \mr{आणि} \mr{थंडी} \mr{पडली}.
\end{quote}

\begin{quote}
\emph{sūrya māvaḷlā āṇi thaṇḍī paḍlī.}
\end{quote}

\begin{quote}
The sun set and it turned cold.
\end{quote}

Past (he …): \mr{तो} \mr{मावळला}, \emph{to māvaḷlā}.

Mostly for the sun and moon; the evening is \mr{मावळतीची} \mr{वेळ}.

\begin{cousinweb}
\textbf{\mr{मावळणे}} (\emph{māvaḷṇe}) is a Marathi verb for the sun or moon going down; its origin is not settled.
\end{cousinweb}

\begin{grammarlens}[title={What you've built}]
One more word for getting around.
\end{grammarlens}

\subsection*{Your turn}
\begin{itemize}
\raggedright
  \item \emph{Say it:} \emph{māvaḷṇe}
  \item \emph{Say it:} \emph{māvaḷṇe}, once more
  \item \emph{Say it:} \emph{sūrya māvaḷlā āṇi thaṇḍī paḍlī.} — The sun set and it turned cold.
  \item \emph{Recall:} say the Marathi for moonlight, then the Marathi for to rise, then say \emph{māvaḷṇe} again
\end{itemize}

\begin{tcolorbox}[breakable,colback=teal!4,colframe=teal!35!black,title={Before you move on}]
What does \mr{मावळणे} mean? (\textquotedblleft{}To set\textquotedblright{}.) Say it once more.
\end{tcolorbox}
\section[\texorpdfstring{\mr{चमकणे} (camakṇe)}{चमकणे (camakṇe)}]{\mr{चमकणे} (camakṇe) — to shine; to flash}

\label{lesson:MR-C441-camakne}



\begin{quote}
Before the new one: say the Marathi for to rise, then the Marathi for to set.
\end{quote}

\subsection*{You'll want to know: \mr{चमकणे}}
\textbf{\mr{चमकणे}} — \emph{camakṇe} — \textquotedblleft{}to shine; to flash\textquotedblright{}.

\begin{quote}
\mr{रात्री} \mr{आकाशात} \mr{तारे} \mr{चमकतात}.
\end{quote}

\begin{quote}
\emph{rātrī ākāśāt tāre camaktāt.}
\end{quote}

\begin{quote}
At night the stars shine in the sky.
\end{quote}

Past (he …): \mr{तो} \mr{चमकला}, \emph{to camaklā}.

A short bright shine: \mr{वीज} \mr{चमकली}, lightning flashed.

\begin{cousinweb}
\textbf{\mr{चमकणे}} (\emph{camakṇe}) is inherited: CDIAL lists Marathi \emph{ċamakṇẽ}, \textquotedblleft{}to glitter\textquotedblright{}, under \emph{*camakka}, \textquotedblleft{}sudden movement\textquotedblright{} (CDIAL 4676).
\end{cousinweb}

\begin{grammarlens}[title={What you've built}]
One more word for getting around.
\end{grammarlens}

\subsection*{Your turn}
\begin{itemize}
\raggedright
  \item \emph{Say it:} \emph{camakṇe}
  \item \emph{Say it:} \emph{camakṇe}, once more
  \item \emph{Say it:} \emph{rātrī ākāśāt tāre camaktāt.} — At night the stars shine in the sky.
  \item \emph{Recall:} say the Marathi for to rise, then the Marathi for to set, then say \emph{camakṇe} again
\end{itemize}

\begin{tcolorbox}[breakable,colback=teal!4,colframe=teal!35!black,title={Before you move on}]
What does \mr{चमकणे} mean? (\textquotedblleft{}To shine; to flash\textquotedblright{}.) Say it once more.
\end{tcolorbox}
